\section*{Chapter 6 --- The interface and the stub}
\addcontentsline{toc}{section}{Chapter 6 --- The interface and the stub}

\solution{ex:6:1}{Predict the compiler}
\ans{a} It compiles without the slightest objection, and that is the whole problem. Without
\code{virtual}, which destructor is called is decided by the pointer's \emph{static} type, that
is, \code{Interface}. \code{Stub}'s destructor never runs, its members are never destroyed, and
the standard calls that undefined behaviour. For \code{Stub} it is not noticed --- it owns nothing
--- but for a class holding a resource it leaks, silently, on every deletion. No test catches it.
\ans{b} Without the attribute, \code{driver.send(frame);} is a perfectly legal statement the
compiler has no reason to complain about. With \code{[[nodiscard]]} it becomes a warning, and with
\code{-Werror} an error. That turns the most common driver bug --- believing that a frame went out
when the hardware was busy --- from a debugging session into a compilation that stops.
\ans{c} With \code{override} the misspelling is a compile error: there is no virtual method by
that name to override. Without \code{override} it compiles, and you have instead \emph{added} a
new method. The base class's pure virtual method is then still unimplemented, so the class remains
abstract --- and the error message you get is about the type not being instantiable, far from the
misspelled line. That is why \code{override} is mandatory.

\solution{ex:6:2}{One method too many}
\ans{a} Nothing meaningful. The stub has no registers; it would have to return an invented
constant, that is, to lie.
\ans{b} \code{EchoNode} becomes tied to the hardware \emph{having} registers. The test for whether
a method belongs in the interface is exactly this: can both implementations fulfil it honestly? If
they cannot, it does not belong there.
\ans{c} Every caller has to change, since the meaning of the bits is part of the register map and
not of the interface. The interface would then inherit the hardware's versions, which is precisely
what it exists to avoid.
\ans{d} In \code{Spi}, as a private helper function --- or, for debugging, in a separate
\code{PrintTransport} below the \code{ByteTransport} seam, where it is invisible to anybody above.

\solution{ex:6:3}{The polymorphism}
The exercise does not ask for code; the questions are answered here.
\ans{a} No. The function takes a \code{driver::can::Interface\&} and uses only the interface's four
methods.
\ans{b} Because it can then be run against any implementation at all. Written against \code{Stub}
it would have been a test of \code{Stub}; written against the interface it is a test of the
\emph{contract}, and the contract is what both \code{Stub} and \code{Spi} promise.
\ans{c} That \code{Spi} really does fulfil the same contract: the same return values in the same
situations. Above all, \code{send()} has to return \code{false} when the hardware is busy and
\code{receive()} has to return \code{false} when there is no frame --- otherwise the function
passes for the stub and fails for the driver.

\solution{ex:6:4}{The limits of the stub}
\ans{a} Errors (\code{hasError()} is always \code{false}), busy hardware (\code{send()} always
succeeds), a queue (a second \code{send()} overwrites the first), and everything CAN does:
arbitration, CRC, bit stuffing.
\ans{b} Busy hardware, for example: a counter that makes the next \code{send()} return
\code{false}, set by a test method. That is the simplest, since it requires no new state beyond a
flag.
\ans{c} Every simulated behaviour is code that can be wrong. A stub that simulates a queue, errors
and timing is a second implementation of the hardware --- and then you no longer know whether a
failing test is due to the driver or to the stub. The limit lies roughly where the stub would need
a test case of its own: if it does, the behaviour belongs in a more realistic double, such as
\code{BankTransport} in \kapref{ch:testning}.
